I wish to express my deepest gratitude to Yushuo Xiao,
first for proposing this thesis project and for accepting me as his student.
Working on a project involving program verification and interactive theorem provers has been a true delight.
I always looked forward to our weekly meetings.
I am grateful for the chance to share my progress and for receiving valuable direction, feedback, and encouragement.
His help with the technical aspects of the project,
especially debugging and navigating Isabelle/HOL, was indispensable,
and I would not have made it far without his guidance.

I would also like to extend a thank you to Prof. Dr. Peter Müller for allowing me to work on such an interesting project
within his research group and for his insightful questions during my presentations.
This project is built upon the prior work of Dr. Gaurav Parthasarathy,
which was supervised by Prof. Müller, whose courses,
such as Program Verification, initially sparked my interest in the field.

Additionally, I would like to thank Prof. Dr. Ralf Jung for teaching the course Formal Foundations of Programming Languages.
His instruction in the formal semantics of programming languages and the interactive theorem prover Coq (now Rocq) helped me find my footing much more quickly in this project.
I also wish to thank Kevin Buzzard and Mohammad Pedramfar for developing "The Natural Number Game,"
which introduced me to Lean and to the world of theorem provers.
I am also grateful to "CADMO - Center for Algorithms, Discrete Mathematics and Optimization" for providing the publicly available \LaTeX\ thesis template\footnote{\url{https://cadmo.ethz.ch/education/thesis/template.html}}.

I am deeply grateful to my friends: Ivana Klasovitá, for her long-time friendship 
and for always being there to talk to,
Remo Spichtig, for the consistently interesting conversations and for the invitations to his board game evenings,
and Kwok Wai Lui, for the exchange of snacks and rants while we worked side-by-side on our respective theses.

Furthermore, I want to thank the sport of volleyball and everyone who enables me to play, including the ASVZ and VBC Rämi.
Engaging in a sport that is physically exerting, sociable, and focused on teamwork provided a vital and healthy balance to my academic work.
I also thank the Swiss Informatics and Mathematical Olympiad for deepening my interest in these fields and for providing me the opportunity to volunteer and give back to the community.

Last but not least, I want to thank my family for their encouragement throughout my studies and for their generous financial support.